We ask whether a one-layer transformer learns $(a+b)\bmod 23$ faster when training pairs are presented in an operand-size curriculum, in the reverse order, or uniformly, and how weight decay and training fraction change the answer. On CPU, with a 4000-step budget, 5 seeds at the main cell (weight decay 1, training fraction 0.5) and 3 seeds in a weight-decay by training-fraction grid, the curriculum did not help: mean steps to 95\% held-out accuracy (censored at 4000) were 2945 for uniform sampling, 3955 for the curriculum and 3990 for the anti-curriculum, with 3, 1 and 1 of 5 seeds reaching 95\%. The curriculum--uniform difference is not significant at 0.05 (Welch and paired tests, 5 seeds). Two further seeds run afterwards (5 and 6) went the other way for the curriculum (uniform 0 of 2, curriculum 1 of 2 reaching 95\%); pooled over 7 seeds the means are \rhval{pooled7/uniform-sampling/wd1_f0.5/steps_to_95/mean} (uniform) and \rhval{pooled7/operand-size-curriculum/wd1_f0.5/steps_to_95/mean} (curriculum), again not significant, so no ordering showed a benefit and a loss is not established. Within the curriculum family, a pacing sweep shows that a short curriculum ($T_c=250$) reached 95\% in all 5 seeds where $T_c=1000$ reached it in 1 of 5; a shuffled-order control with the same pacing did no better than the operand-size order. With zero weight decay no run in any condition generalised within the budget, and in 5 of 9 grid cells no run of any system did, so delayed generalisation was only observed in part of the grid. The sign of the curriculum effect varied across the two cells where curriculum and uniform both reached 95\% in all seeds, by amounts well inside seed noise. This is a single-modulus, single-architecture, few-seed study; it supports ``no benefit found'', not ``no benefit exists''.
